% Options for packages loaded elsewhere
\PassOptionsToPackage{unicode}{hyperref}
\PassOptionsToPackage{hyphens}{url}
\PassOptionsToPackage{dvipsnames,svgnames,x11names}{xcolor}
\documentclass[
  11pt,
  a4paper,
]{article}
\usepackage{xcolor}
\usepackage[margin=1in]{geometry}
\usepackage{amsmath,amssymb}
\setcounter{secnumdepth}{-\maxdimen} % remove section numbering
\usepackage{iftex}
\ifPDFTeX
  \usepackage[T1]{fontenc}
  \usepackage[utf8]{inputenc}
  \usepackage{textcomp} % provide euro and other symbols
\else % if luatex or xetex
  \usepackage{unicode-math} % this also loads fontspec
  \defaultfontfeatures{Scale=MatchLowercase}
  \defaultfontfeatures[\rmfamily]{Ligatures=TeX,Scale=1}
\fi
\usepackage{lmodern}
\ifPDFTeX\else
  % xetex/luatex font selection
  \ifLuaTeX
    \usepackage{luaotfload}
    \directlua{luaotfload.add_fallback("mainfontfallback",{
      "NotoSerifBengali:mode=harf"
    })}
  \fi
  \setmainfont[,RawFeature={fallback=mainfontfallback}]{Latin Modern Roman}
  \ifLuaTeX
    \usepackage{luaotfload}
    \directlua{luaotfload.add_fallback("sansfontfallback",{
      "NotoSansBengali:mode=harf"
    })}
  \fi
  \setsansfont[,RawFeature={fallback=sansfontfallback}]{Latin Modern Sans}
  \ifLuaTeX
    \usepackage{luaotfload}
    \directlua{luaotfload.add_fallback("monofontfallback",{
      "NotoSansBengali:mode=harf"
    })}
  \fi
  \setmonofont[,RawFeature={fallback=monofontfallback}]{DejaVu Sans Mono}
\fi
% Use upquote if available, for straight quotes in verbatim environments
\IfFileExists{upquote.sty}{\usepackage{upquote}}{}
\IfFileExists{microtype.sty}{% use microtype if available
  \usepackage[]{microtype}
  \UseMicrotypeSet[protrusion]{basicmath} % disable protrusion for tt fonts
}{}
\makeatletter
\@ifundefined{KOMAClassName}{% if non-KOMA class
  \IfFileExists{parskip.sty}{%
    \usepackage{parskip}
  }{% else
    \setlength{\parindent}{0pt}
    \setlength{\parskip}{6pt plus 2pt minus 1pt}}
}{% if KOMA class
  \KOMAoptions{parskip=half}}
\makeatother
\usepackage{longtable,booktabs,array}
\usepackage{calc} % for calculating minipage widths
% Correct order of tables after \paragraph or \subparagraph
\usepackage{etoolbox}
\makeatletter
\patchcmd\longtable{\par}{\if@noskipsec\mbox{}\fi\par}{}{}
\makeatother
% Allow footnotes in longtable head/foot
\IfFileExists{footnotehyper.sty}{\usepackage{footnotehyper}}{\usepackage{footnote}}
\makesavenoteenv{longtable}
\setlength{\emergencystretch}{3em} % prevent overfull lines
\providecommand{\tightlist}{%
  \setlength{\itemsep}{0pt}\setlength{\parskip}{0pt}}
% paper_preamble.tex — shared LaTeX preamble for the PCMA/SI paper set.
%
% Compiled with lualatex (NOT pdflatex) because several papers contain
% Bengali script: the CLL paper realises one cognitive state into
% English, Bengali and Banglish, so the Bengali output is EVIDENCE and
% cannot be transliterated away without destroying the point being made.
%
% Font selection is NOT done here. pandoc >= 3 has native
% `mainfontfallback` / `monofontfallback` variables that emit the
% luaotfload fallback registration in the correct position relative to
% unicode-math; doing it by hand in this header lands after unicode-math
% and breaks font loading outright. See tools/build_papers.sh.

\usepackage{amsmath}
\usepackage{amssymb}
\usepackage{mathtools}   % \xleftrightarrow etc. used in cll.md §5
\usepackage{booktabs}

% ---------------------------------------------------------------------
% Unicode symbols appearing in PROSE rather than inside $...$.
%
% The papers were written as markdown and read as markdown, so they use
% real characters — "the ω-boost", "α-child", "ΔR coherence" — instead of
% LaTeX macros. Under unicode-math those codepoints are math-mode-only
% and abort the build ("\symbballowed only in math mode"). Mapping each
% to \ensuremath keeps the markdown readable on GitHub AND compilable,
% which is the whole point of keeping one source.
% ---------------------------------------------------------------------
\usepackage{newunicodechar}

% Greek, lower case
\newunicodechar{α}{\ensuremath{\alpha}}
\newunicodechar{β}{\ensuremath{\beta}}
\newunicodechar{γ}{\ensuremath{\gamma}}
\newunicodechar{δ}{\ensuremath{\delta}}
\newunicodechar{ε}{\ensuremath{\varepsilon}}
\newunicodechar{η}{\ensuremath{\eta}}
\newunicodechar{θ}{\ensuremath{\theta}}
\newunicodechar{κ}{\ensuremath{\kappa}}
\newunicodechar{λ}{\ensuremath{\lambda}}
\newunicodechar{μ}{\ensuremath{\mu}}
\newunicodechar{π}{\ensuremath{\pi}}
\newunicodechar{ρ}{\ensuremath{\rho}}
\newunicodechar{σ}{\ensuremath{\sigma}}
\newunicodechar{τ}{\ensuremath{\tau}}
\newunicodechar{φ}{\ensuremath{\varphi}}
\newunicodechar{χ}{\ensuremath{\chi}}
\newunicodechar{ψ}{\ensuremath{\psi}}
\newunicodechar{ω}{\ensuremath{\omega}}
\newunicodechar{ν}{\ensuremath{\nu}}
\newunicodechar{ξ}{\ensuremath{\xi}}
\newunicodechar{ζ}{\ensuremath{\zeta}}
\newunicodechar{ι}{\ensuremath{\iota}}

% Greek, upper case
\newunicodechar{Δ}{\ensuremath{\Delta}}
\newunicodechar{Φ}{\ensuremath{\Phi}}
\newunicodechar{Σ}{\ensuremath{\Sigma}}
\newunicodechar{Ψ}{\ensuremath{\Psi}}
\newunicodechar{Ω}{\ensuremath{\Omega}}
\newunicodechar{Γ}{\ensuremath{\Gamma}}
\newunicodechar{Λ}{\ensuremath{\Lambda}}
\newunicodechar{Θ}{\ensuremath{\Theta}}
\newunicodechar{Π}{\ensuremath{\Pi}}

% Operators and relations
\newunicodechar{×}{\ensuremath{\times}}
\newunicodechar{·}{\ensuremath{\cdot}}
\newunicodechar{±}{\ensuremath{\pm}}
\newunicodechar{−}{\ensuremath{-}}
\newunicodechar{≈}{\ensuremath{\approx}}
\newunicodechar{≤}{\ensuremath{\leq}}
\newunicodechar{≥}{\ensuremath{\geq}}
\newunicodechar{≠}{\ensuremath{\neq}}
\newunicodechar{∈}{\ensuremath{\in}}
\newunicodechar{∂}{\ensuremath{\partial}}
\newunicodechar{∇}{\ensuremath{\nabla}}
\newunicodechar{√}{\ensuremath{\sqrt{}}}
\newunicodechar{∞}{\ensuremath{\infty}}
\newunicodechar{∝}{\ensuremath{\propto}}
\newunicodechar{∑}{\ensuremath{\sum}}
\newunicodechar{∫}{\ensuremath{\int}}
\newunicodechar{ℝ}{\ensuremath{\mathbb{R}}}
\newunicodechar{ℕ}{\ensuremath{\mathbb{N}}}
\newunicodechar{‖}{\ensuremath{\|}}
\newunicodechar{⊕}{\ensuremath{\oplus}}
\newunicodechar{⊗}{\ensuremath{\otimes}}
\newunicodechar{∀}{\ensuremath{\forall}}
\newunicodechar{∃}{\ensuremath{\exists}}
\newunicodechar{¬}{\ensuremath{\neg}}
\newunicodechar{∧}{\ensuremath{\wedge}}
\newunicodechar{∨}{\ensuremath{\vee}}

% Arrows
\newunicodechar{→}{\ensuremath{\rightarrow}}
\newunicodechar{←}{\ensuremath{\leftarrow}}
\newunicodechar{↔}{\ensuremath{\leftrightarrow}}
\newunicodechar{↑}{\ensuremath{\uparrow}}
\newunicodechar{↓}{\ensuremath{\downarrow}}
\newunicodechar{⇒}{\ensuremath{\Rightarrow}}
\newunicodechar{⇐}{\ensuremath{\Leftarrow}}
\newunicodechar{⇔}{\ensuremath{\Leftrightarrow}}

% Super/subscripts written as literal characters
\newunicodechar{²}{\ensuremath{^{2}}}
\newunicodechar{³}{\ensuremath{^{3}}}
\newunicodechar{⁴}{\ensuremath{^{4}}}
\newunicodechar{⁵}{\ensuremath{^{5}}}
\newunicodechar{⁶}{\ensuremath{^{6}}}
\newunicodechar{₀}{\ensuremath{_{0}}}
\newunicodechar{₁}{\ensuremath{_{1}}}
\newunicodechar{₂}{\ensuremath{_{2}}}
\newunicodechar{₃}{\ensuremath{_{3}}}
\newunicodechar{ᵀ}{\ensuremath{^{\mathsf{T}}}}

% Status glyphs. The papers use emoji as status markers in tables
% (implemented / failed / warning). Mapped to typographic equivalents:
% a colour emoji font in lualatex is unreliable, and the MEANING is what
% carries — a reader needs to see "this one failed", not a red cross.
\newunicodechar{✅}{\ensuremath{\checkmark}}
\newunicodechar{✓}{\ensuremath{\checkmark}}
\newunicodechar{❌}{\ensuremath{\times}}
\newunicodechar{✗}{\ensuremath{\times}}
\newunicodechar{⚠}{\textbf{!}}
\newunicodechar{🧩}{\textbf{[piece]}}
\newunicodechar{📋}{\textbf{[list]}}
\newunicodechar{🟡}{\textbf{[partial]}}
\newunicodechar{🔴}{\textbf{[dead]}}
\newunicodechar{🟢}{\textbf{[live]}}
\newunicodechar{▼}{\ensuremath{\blacktriangledown}}
\newunicodechar{►}{\ensuremath{\blacktriangleright}}
\newunicodechar{□}{\ensuremath{\square}}
\newunicodechar{█}{\rule{0.6em}{0.9ex}}
\newunicodechar{•}{\textbullet}
\newunicodechar{†}{\textdagger}
\newunicodechar{‑}{-}
\newunicodechar{‐}{-}
\newunicodechar{…}{\ldots}
\newunicodechar{।}{\textbar}   % danda: rendered by the Bengali fallback in text
\newunicodechar{⁻}{\ensuremath{^{-}}}

% Zero-width formatting codepoints. These are invisible in an editor but
% are real characters to LaTeX: an emoji written as "⚠️" is TWO
% codepoints (U+26A0 U+FE0F), so declaring only the visible one leaves a
% variation selector with no definition and the build dies on a symbol
% nobody can see. Both are mapped to nothing.
\newunicodechar{^^^^fe0f}{}    % VARIATION SELECTOR-16
\newunicodechar{^^^^200d}{}    % ZERO WIDTH JOINER
\newunicodechar{^^^^200b}{}    % ZERO WIDTH SPACE

\usepackage{longtable}
\usepackage{array}

% Comparison matrices and cost ledgers are wide. Let them shrink to the
% text block rather than overflow the margin.
\usepackage{adjustbox}

% Code blocks: wrap long lines instead of running off the page. Several
% papers quote real C++ with long identifiers.
\usepackage{fvextra}
\fvset{breaklines=true,breakanywhere=true,fontsize=\small}

\setlength{\parskip}{0.5em}
\setlength{\parindent}{0pt}

\usepackage{fancyhdr}
\pagestyle{fancy}
\fancyhf{}
\fancyfoot[C]{\thepage}
\fancyhead[L]{\small\itshape Synthetic Intelligence — research lineage}
\fancyhead[R]{\small\itshape irfan@furylogic.com}
\renewcommand{\headrulewidth}{0.2pt}
\setlength{\headheight}{14pt}
\setlength{\headsep}{18pt}
\usepackage{bookmark}
\IfFileExists{xurl.sty}{\usepackage{xurl}}{} % add URL line breaks if available
\urlstyle{same}
\hypersetup{
  pdftitle={The L9-Cognitive Graph: Nine Layers from Lexical Decomposition to Paradigm Generation},
  pdfauthor={Irfan Mahir --- Independent Researcher, Dhaka, Bangladesh --- irfan@furylogic.com},
  colorlinks=true,
  linkcolor={black},
  filecolor={Maroon},
  citecolor={Blue},
  urlcolor={[HTML]{0B5394}},
  pdfcreator={LaTeX via pandoc}}

\title{The L9-Cognitive Graph: Nine Layers from Lexical Decomposition to
Paradigm Generation}
\author{Irfan Mahir --- Independent Researcher, Dhaka, Bangladesh ---
irfan@furylogic.com}
\date{}

\begin{document}
\maketitle

{
\hypersetup{linkcolor=}
\setcounter{tocdepth}{3}
\tableofcontents
}
\section{The L9-Cognitive Graph: Nine Layers from Lexical Decomposition
to Paradigm
Generation}\label{the-l9-cognitive-graph-nine-layers-from-lexical-decomposition-to-paradigm-generation}

\textbf{Irfan Mahir} \emph{Independent Researcher, Dhaka, Bangladesh}
irfan@furylogic.com

\begin{quote}
\emph{``A graph that answers
\texttt{{[}Apple{]}\ -\/-is\_a-\/-\textgreater{}\ {[}Fruit{]}} is a
dictionary. Maxwell did not consult a dictionary.''}
\end{quote}

\begin{center}\rule{0.5\linewidth}{0.5pt}\end{center}

\subsection{Abstract}\label{abstract}

Probabilistic next-symbol prediction confines a system to interpolating
within existing knowledge. We propose the \textbf{L9-Cognitive Graph}, a
neuro-symbolic architecture organised as nine discrete layers of
traversal logic over a deterministic symbolic graph, ascending from
lexical decomposition to the generation of new axiomatic roots.

The organising claim is that breakthrough reasoning is not semantic
retrieval but \textbf{cross-domain isomorphism}: recognising that the
mathematical structure of one domain describes the unresolved anomalies
of another. Achieving this requires the engine to move from
\emph{constraint enforcement} to \emph{abductive constraint generation}.

This paper specifies all nine layers --- each with its input, operation,
output and failure mode --- rather than listing them. Section 5 states,
layer by layer, which are implemented, which are partial, and which
remain proposals. Four are live, two are partial, and three are not
built.

\begin{center}\rule{0.5\linewidth}{0.5pt}\end{center}

\subsection{1. The Limits of Semantic
Retrieval}\label{the-limits-of-semantic-retrieval}

A graph database mapping
\texttt{{[}Apple{]}\ -\/-is\_a-\/-\textgreater{}\ {[}Fruit{]}} is a
semantic dictionary: it returns what was stored, structured. Human
breakthrough logic does not work this way. Maxwell's unification of
electricity and magnetism was not a lookup; it was the recognition that
one domain's mathematical structure resolved another's anomalies.

Reproducing that requires two capabilities a dictionary lacks. The
engine must be able to notice that its own knowledge is
\textbf{structurally inconsistent} rather than merely incomplete, and it
must be able to \textbf{propose structure that was never asserted} in
response.

The layers below are ordered by how far each departs from retrieval.

\begin{center}\rule{0.5\linewidth}{0.5pt}\end{center}

\subsection{2. Biological Motivation: Two
Topologies}\label{biological-motivation-two-topologies}

The architecture borrows two structures rather than a metaphor.

\textbf{Neocortical columns --- the local graph.} The brain processes
localised, highly deterministic rules: syntax, basic physics. These
correspond to densely connected subgraph clusters with near-unit edge
weights, where traversal is reliable and cheap.

\textbf{The Default Mode Network --- the abductive engine.} The DMN
activates during undirected thought and deep problem-solving, searching
for low-probability, long-distance connections across disparate columns.
It corresponds to an asynchronous background process running
bridge-finding across structurally similar but semantically distant
subgraphs.

The functional distinction that matters: \textbf{the local graph
answers, the DMN wanders.} A system with only the first can never
surprise itself.

\begin{center}\rule{0.5\linewidth}{0.5pt}\end{center}

\subsection{3. The Nine Layers}\label{the-nine-layers}

Each layer is specified by what it takes, what it does, what it emits,
and how it fails. Layers are cumulative: L\(n\) presupposes
L\(1..n{-}1\).

\subsubsection{L1 --- Lexical / Morphological
Decomposition}\label{l1-lexical-morphological-decomposition}

\textbf{Input.} A token that resolves to no node. \textbf{Operation.}
Fracture the string into sub-lexical fragments --- morphemes, syllables,
character clusters --- and project each against the ontology for partial
matches. \texttt{biology} decomposes to \texttt{bio} + \texttt{logy};
\texttt{sudo} to \texttt{su} + \texttt{do}. \textbf{Output.} A set of
candidate bindings for an otherwise orphan node. \textbf{Failure mode.}
Spurious decomposition. Short fragments match everything, so a fragment
carries usable evidence only above a minimum length --- an information
floor, not a stopword list.

This layer exists because a node of degree zero is unreachable, and an
unreachable node makes every higher layer inapplicable.

\subsubsection{L2 --- Ontological
Retrieval}\label{l2-ontological-retrieval}

\textbf{Input.} A resolved node. \textbf{Operation.} Traverse explicit
edges. \texttt{Darwin\ -\textgreater{}\ Human}. \textbf{Output.}
Asserted neighbours. \textbf{Failure mode.} Silence when no edge exists.
This is the honest failure and is preferable to a fabricated one.

L1 and L2 together constitute the dictionary. Everything above is the
argument of this paper.

\subsubsection{L3 --- Transitive
Inference}\label{l3-transitive-inference}

\textbf{Input.} Two chained edges \(A \to B\), \(B \to C\).
\textbf{Operation.} Instantiate the implied edge \(A \to C\), recording
both premises. \textbf{Output.} A derived edge, distinguishable from an
asserted one. \textbf{Failure mode.} Unbounded expansion. Transitive
closure over a large graph is quadratic, so derivation must be bounded
by an explicit budget.

The provenance requirement is not decoration. A derived edge that cannot
name its premises cannot be retracted when one of them is withdrawn.

\subsubsection{L4 --- Causal Abduction}\label{l4-causal-abduction}

\textbf{Input.} An observed end state \(Y\). \textbf{Operation.} Search
\emph{backwards} for a cause \(X\) whose consequences include \(Y\).
This inverts L3's direction of travel. \textbf{Output.} A hypothesis,
not a conclusion. \textbf{Failure mode.} Affirming the consequent. Many
causes produce the same effect, so an abduced \(X\) is a candidate
requiring independent support --- and labelling it otherwise is the
single easiest way for this architecture to produce confident nonsense.

\subsubsection{L5 --- Counterfactual
Simulation}\label{l5-counterfactual-simulation}

\textbf{Input.} A core axiom. \textbf{Operation.} Temporarily invert it
--- \emph{what if gravity repelled?} --- and propagate consequences,
checking whether the resulting system remains consistent.
\textbf{Output.} A stability verdict on the inverted world.
\textbf{Failure mode.} Combinatorial explosion, and contamination of the
base graph if the inversion is not perfectly isolated and reverted.

L5 is how the engine determines which axioms are \emph{load-bearing}: an
axiom whose inversion changes nothing was doing no work.

\subsubsection{L6 --- Analogical Mapping}\label{l6-analogical-mapping}

\textbf{Input.} Two subgraphs from different domains.
\textbf{Operation.} Attempt a structure-preserving \(1{:}1\)
correspondence --- fluid dynamics onto electrical circuits, pressure
\(\to\) voltage, flow \(\to\) current. \textbf{Output.} A mapping, plus
the edges of one domain transferred onto the other as hypotheses.
\textbf{Failure mode.} Superficial isomorphism. Two subgraphs can share
shape while the correspondence means nothing, so a mapping must be
scored by more than topology.

\textbf{This is the first layer of genuine novelty}, because its output
is not entailed by either domain alone.

\subsubsection{L7 --- Anomaly Detection}\label{l7-anomaly-detection}

\textbf{Input.} The graph. \textbf{Operation.} Scan for nodes of high
connectivity carrying \emph{contradictory} constraints --- a concept
that participates widely while holding incompatible commitments. This is
the \textbf{Maxwell gap}: the signature of a theory that is incomplete
at a specific, locatable point. \textbf{Output.} A ranked list of
structural contradictions. \textbf{Failure mode.} Mistaking noise for
anomaly. Extraction error produces contradictions that indicate a bad
parse rather than a gap in the world.

L7 is the layer that lets the system be \emph{dissatisfied} with what it
knows, which is the precondition for L8 and L9.

\subsubsection{L8 --- Synthesis (The DMN
Layer)}\label{l8-synthesis-the-dmn-layer}

\textbf{Input.} An L7 anomaly. \textbf{Operation.} Generate a
\textbf{bridge node} resolving it, constrained to satisfy both of the
subgraphs it joins. \textbf{Output.} A new vertex that appeared in no
source. \textbf{Failure mode.} Unfalsifiable patching. A bridge invented
purely to absorb a contradiction, with no independent consequences,
explains nothing --- it renames the problem.

The discipline required here is that a synthesised node must make
predictions beyond the anomaly that motivated it.

\subsubsection{L9 --- Paradigm Generation}\label{l9-paradigm-generation}

\textbf{Input.} A region of the graph that is verbose, irregular and
patched. \textbf{Operation.} Define a new \textbf{axiomatic root} from
which the messy region follows, collapsing many special cases into fewer
general ones. \textbf{Output.} A reorganisation of the knowledge base,
not an addition to it. \textbf{Failure mode.} Losing content. A
simplification that fails to entail what it replaced is not a paradigm;
it is deletion.

L9 is the only layer whose success criterion is \emph{compression}: the
same commitments held in less structure.

\begin{center}\rule{0.5\linewidth}{0.5pt}\end{center}

\subsection{4. Scoring Novelty}\label{scoring-novelty}

Novelty is not randomness --- it is the discovery of hidden structural
isomorphism. Given subgraphs \(A\) and \(B\), we seek pairs that are
\textbf{semantically distant} but \textbf{structurally identical}:

\[
\Phi(A, B) \;=\; \frac{I(A, B)}{e^{-k \cdot d(A, B)}}
\]

where \(I(A,B)\) is a maximal common subgraph isomorphism score,
\(d(A,B)\) is the shortest path between the subgraph roots, and \(k\) is
a scaling constant.

The exponential in the denominator is the load-bearing choice: it makes
\(\Phi\) grow with distance, so a structural match between
\emph{adjacent} concepts scores poorly. Two neighbouring nodes sharing
shape is expected. Two nodes twelve hops apart sharing shape is a
candidate breakthrough. Pairs exceeding a threshold \(\tau\) are
flagged.

\textbf{A caution this formulation invites.} \(\Phi\) rewards distance
unboundedly, so with a finite isomorphism score the highest-scoring
pairs will tend to be the most remote --- which includes pairs that are
remote because one of them is extraction noise. Any deployment needs a
floor on subgraph quality before distance is rewarded.

\begin{center}\rule{0.5\linewidth}{0.5pt}\end{center}

\subsection{5. Implementation Status}\label{implementation-status}

The paper's claims are separated from the engine's behaviour here, layer
by layer.

\begin{longtable}[]{@{}
  >{\raggedright\arraybackslash}p{(\linewidth - 4\tabcolsep) * \real{0.3333}}
  >{\raggedright\arraybackslash}p{(\linewidth - 4\tabcolsep) * \real{0.3333}}
  >{\raggedright\arraybackslash}p{(\linewidth - 4\tabcolsep) * \real{0.3333}}@{}}
\toprule\noalign{}
\begin{minipage}[b]{\linewidth}\raggedright
Layer
\end{minipage} & \begin{minipage}[b]{\linewidth}\raggedright
Status
\end{minipage} & \begin{minipage}[b]{\linewidth}\raggedright
Where, or what is missing
\end{minipage} \\
\midrule\noalign{}
\endhead
\bottomrule\noalign{}
\endlastfoot
\textbf{L1} Morphological & \textbf{live} & bigram sensory lanes; graded
recognition survives misspelling, 7/7 measured \\
\textbf{L2} Retrieval & \textbf{live} & typed concept graph traversal \\
\textbf{L3} Transitive & \textbf{live} & forward chaining with premise
citation; derived edges name both parents \\
\textbf{L4} Abduction & \textbf{partial} & abduction over upstream
cavities exists; no general backward search from an arbitrary \(Y\) \\
\textbf{L5} Counterfactual & \textbf{not built} & requires isolated
graph forking with guaranteed revert; no such isolation exists \\
\textbf{L6} Analogical & \textbf{live} & shared-target intersection plus
rubber-tunnel bridging, with 20-seed reproducibility data; the Fisher
metric supplies a principled distance \\
\textbf{L7} Anomaly & \textbf{partial} & contradiction detection and
isolated-subgraph/logic-gap scans exist; no ranked Maxwell-gap
register \\
\textbf{L8} Synthesis & \textbf{partial} & fission spawns new vertices
to resolve contradiction, and mutation forces distant-pair hypotheses
that decay if unsupported --- but neither is \emph{driven by} an L7
anomaly report \\
\textbf{L9} Paradigm & \textbf{not built} & no mechanism reorganises the
graph around a new root; macro-rule induction compresses
\emph{derivations}, not \emph{axioms}, which is a weaker relative \\
\end{longtable}

\textbf{The honest summary.} The lower half of the stack is real. The
upper half is partially present but \emph{disconnected}: the engine has
anomaly detection and it has synthesis, but synthesis does not consume
anomaly reports. Wiring L7 into L8 is the single highest-value unbuilt
link, and it is engineering rather than research.

L5 and L9 are genuinely open. L9 in particular is not a matter of
effort: no criterion is currently specified for when a proposed
axiomatic root is \emph{better} than the structure it would replace, and
without that criterion the layer cannot be evaluated even if built.

\begin{center}\rule{0.5\linewidth}{0.5pt}\end{center}

\subsection{6. Limitations}\label{limitations}

\textbf{The layer ordering is asserted.} Nine layers with clean
boundaries is a presentational convenience. In the running engine, L1
recognition and L6 bridging interact directly, and L3 derivation feeds
L7 detection without passing through L4 or L5. The stack is a taxonomy,
not a pipeline.

\textbf{\(\Phi\) is untuned.} The constant \(k\) and threshold \(\tau\)
have no principled setting, and §4 notes the formulation's bias toward
remote noise. The Fisher-information distance is a more defensible
replacement for \(d\), and supersedes the shortest-path form where both
are available.

\textbf{No layer above L3 has an oracle.} L3's output is checkable by
construction. An L6 mapping, an L8 bridge or an L9 paradigm has no
external verdict available --- which is precisely the gap the
compile-verify loop closes for code, and does not close for physics or
biology.

\begin{center}\rule{0.5\linewidth}{0.5pt}\end{center}

\subsection{7. Conclusion}\label{conclusion}

The L9 stack describes the ascent from a semantic dictionary to an
isomorphic constraint-satisfaction engine --- one able to observe
structural contradictions in what it holds and calculate the structure
required to unify them.

Four layers are live and two are partial. The most useful outcome of
writing the stack out in full is not the top of it but the middle: it
makes visible that anomaly detection and synthesis both exist and are
\textbf{not connected to each other}, which was not apparent while the
layers were a list of one-line descriptions.

\begin{center}\rule{0.5\linewidth}{0.5pt}\end{center}

\subsection{References}\label{references}

\begin{itemize}
\tightlist
\item
  Mahir, I. \emph{From Artificial to Synthetic Intelligence.} §4.4 ---
  the nine reasoning primitives implementing L2--L6.
\item
  Mahir, I. \emph{RB+G: Gated Recursive Bifurcation with Analogical
  Bridging.} --- L6 and part of L8, with reproducibility data.
\item
  Mahir, I. \emph{Conceptual Mutation Architecture (CMA).} --- forced
  distant-pair hypothesis with decay; the other half of L8.
\item
  Mahir, I. \emph{Morphological Abduction Architecture (MAA).} --- L1
  and L4.
\item
  Mahir, I. \emph{A Boltzmann Fisher-Information Metric for Graph-Valued
  Cognitive States.} --- a principled replacement for \(d(A,B)\) in §4.
\item
  Mahir, I. \emph{JAS: Einstein's Epistemology as a Running
  Architecture.} --- the oracle problem of §6.
\item
  Implementation: \texttt{src/layer2.hpp}, \texttt{src/analogy.hpp},
  \texttt{src/inference.hpp}, \texttt{src/bisociation.hpp}.
\end{itemize}

\end{document}
